\section{Estimators, Aggregation, and Inferential Scope}
\label{app:inference}

Section~\ref{sec:estimands} states the design-level estimands and inference units.
This appendix answers four narrower questions: how the predictive quantities are related, how fitting and aggregation are nested, how uncertainty and decision rules are assigned, and where information-theoretic analogies stop.
Appendix~\ref{app:data-targets} owns the inputs and preprocessing definitions, Appendix~\ref{app:training} owns the training objectives and retained model states, and Section~\ref{sec:results} owns the empirical outcomes.
The quantities below are assay-conditional predictive contrasts.
They are neither causal effects nor direct estimates of information-theoretic objects.

\subsection{Predictive estimands and ridge estimation}

For response coordinate \(j\), Section~\ref{sec:measurement-prerequisites} defines the held-out score \(R_j^2(M)\) for a readout design \(M\) in Equation~\ref{eq:heldout-r2}.
Let \(N\) denote the declared nuisance design and \(F=[N,X]\) the full design after adding contextual \ac{lm} representations \(X\).
The nuisance-only and full scores are \(R_j^2(N)\) and \(R_j^2(F)\), and the reported semipartial increment is
\[
  \Delta R_{\mathrm{semi},j}^2(X)
  =
  R_j^2(F)-R_j^2(N).
\]
This difference asks how much held-out predictive value is gained by adding \(X\) to the nuisance design.
Because the nuisance-only and full ridge readouts are separately regularized and evaluated out of sample, the difference can be negative.
It is not a nonnegative variance decomposition. \evd{E002}\evd{E006}

Population partial \(R^2\) answers a different question.
Let \(\mathcal R_{N,j}^2\) and \(\mathcal R_{F,j}^2\) be the population coefficients of determination for nested, unregularized linear projections on \(N\) and \(F\).
When \(\mathcal R_{N,j}^2<1\),
\begin{equation}
  R_{\mathrm{partial},j}^2
  =
  \frac{\mathcal R_{F,j}^2-\mathcal R_{N,j}^2}
       {1-\mathcal R_{N,j}^2}
  =
  1-
  \frac{\operatorname{Var}\!\left(Y_j-\Pi_FY_j\right)}
       {\operatorname{Var}\!\left(Y_j-\Pi_NY_j\right)},
  \label{eq:partial-r2}
\end{equation}
where \(\Pi_N\) and \(\Pi_F\) are the corresponding population linear projections.
Equation~\ref{eq:partial-r2} measures the fraction of variance left after the nuisance projection that is removed by adding \(X\).
The semipartial quantity instead measures the gain relative to total outcome variance.
The thesis reports the cross-validated semipartial increment, not a transformed cross-validated partial coefficient.

Aggregation across target or response coordinates defines another distinction.
For \(J\) held-out coordinates, let \(\mathrm{SSE}_j\) and \(\mathrm{SST}_j\) denote the coordinate-specific residual and total sums of squares.
The two aggregations used in the thesis are
\begin{equation}
  \overline{R^2}_{\mathrm{equal}}
  =
  \frac{1}{J}\sum_{j=1}^{J}
  \left(1-\frac{\mathrm{SSE}_j}{\mathrm{SST}_j}\right),
  \qquad
  R^2_{\mathrm{pooled}}
  =
  1-\frac{\sum_{j=1}^{J}\mathrm{SSE}_j}
           {\sum_{j=1}^{J}\mathrm{SST}_j}.
  \label{eq:coordinate-aggregation}
\end{equation}
The first expression gives every coordinate equal weight.
The second weights coordinates by their held-out variance, so a nearly constant coordinate cannot contribute as much as a variable coordinate.
The WikiText synthetic-target recovery assay uses equal-coordinate averaging, whereas the saved-student target-retention analysis pools sums of squares across the fixed target coordinates.
Neither calculation treats coordinates as independent repetitions. \evd{E016}\evd{E030}

Ridge regression supplies a regularized linear predictor.
Let \(y\in\mathbb R^n\) be a centered response vector, let \(M\in\mathbb R^{n\times p}\) be the centered and standardized predictor matrix, and let \(\beta\in\mathbb R^p\) be its coefficient vector.
If
\[
  y=M\beta+\varepsilon,
  \qquad
  \varepsilon\sim\mathcal N(0,\sigma^2I),
  \qquad
  \beta\sim\mathcal N(0,\tau_\beta^2I),
\]
then the posterior mode minimizes
\[
  \lVert y-M\beta\rVert_2^2
  +
  \lambda\lVert\beta\rVert_2^2,
  \qquad
  \lambda=\frac{\sigma^2}{\tau_\beta^2}.
\]
This is a maximum a posteriori interpretation only under the stated Gaussian likelihood, isotropic zero-mean prior, and scaling convention.
The fitted linear predictor represents the conditional mean only when that mean is linear and the model is correctly specified.

The sentence-level Tuckute and WikiText readouts use the candidate penalties
\[
\{1,10,100,1000,10000\}.
\]
Each multi-output fit selects one scalar penalty on its training partition by ridge cross-validation; it does not select a separate penalty for every response or target coordinate.
The naturalistic LeBel readouts use
\[
\{10,10^2,10^3,10^4,10^5,10^6,10^7\},
\]
again selecting one multi-output penalty from training stories only.

The comparator-adequacy audit uses
\[
\{0.1,1,10,100,1000\} % \evd{E026}
\]
and selects one penalty by aggregate held-out \(R^2\) across both target families and both fixed coordinate samples within five contiguous inner folds.
These selection rules are nested inside the corresponding training partition; the untouched outer block or held-out corpus is used once for evaluation.\evd{E002}\evd{E006}\evd{E016}\evd{E025}\evd{E026}\evd{E030}
In the executed assays, \(\lambda\) is selected from a finite grid using the relevant training complement; operationally, it is predictive regularization rather than a fixed scientific prior.
Neither this derivation nor a positive held-out score makes the fitted coefficients causal, identifies a biological mechanism, or shows that an \ac{lm} and the brain implement the same computation.

\subsection{Cross-validation, aggregation, and inference units}

The governing leakage rule is simple: every data-adaptive operation that can use stimulus-domain information is fitted on the relevant training complement.
Fixed protocol operations may be applied before the split, but learned centering, scaling, dimensionality reduction, penalty selection, and readout fitting remain inside it.
Appendix~\ref{app:data-targets} specifies the exact fixed preprocessing and split definitions.

The naturalistic assays require one explicit distinction.
Feature resampling, edge slicing, within-story feature and \ac{bold} z-scoring, and finite-impulse-response delay copies are fixed story-local operations.
Because the held-out story's response normalization uses that story's own moments, the resulting estimand is conditional on this normalization protocol.
Across-story \ac{pca}, assembled-design scaling, response centering, ridge selection, and readout fitting are learned only from the readout-training stories.
The 50-component linear target-projection assay has one other declared exception: it reuses target-standardization moments fixed during target construction, while its target \ac{pca} and readouts remain fold-local. \evd{E006}\evd{E013}\evd{E017}\evd{E030}

Table~\ref{tab:cv-nesting} records the held-out boundary, learned operations, and final inference unit for each scientific role.

\begingroup
\footnotesize
\setlength{\tabcolsep}{2pt}
\renewcommand{\arraystretch}{1}
\begin{longtable}{@{}>{\raggedright\arraybackslash}p{27mm}>{\raggedright\arraybackslash}p{29mm}>{\raggedright\arraybackslash}p{47mm}>{\raggedright\arraybackslash}p{31mm}@{}}
  \caption{Cross-validation nesting and inference units by scientific role.}
  \label{tab:cv-nesting}\\
  \toprule
  Scientific role & Held-out boundary & Operations learned inside the boundary & Aggregation and final unit \\
  \midrule
  \endfirsthead
  \multicolumn{4}{c}{\tablename~\thetable{} (continued)} \\
  \toprule
  Scientific role & Held-out boundary & Operations learned inside the boundary & Aggregation and final unit \\
  \midrule
  \endhead
  \midrule
  \multicolumn{4}{r}{Continued on next page} \\
  \endfoot
  \bottomrule
  \endlastfoot
  Controlled regional predictivity assay &
  Five contiguous Tuckute folds, each with an 800-sentence training complement and a 200-sentence test block &
  Centering, scaling, \ac{pca}, one multi-output ridge penalty, and the response readout &
  Declared \acp{roi} and folds are averaged; the sentence set is fixed, and untrained-model seeds are technical repetitions \evd{E002} \\
  \midrule
  Distillation-headroom analysis &
  The same five Tuckute folds for brain predictivity and a fixed WikiText probe for language quality &
  Brain-response transforms and readouts are fitted inside each fold; quality scores use the fixed scoring protocol &
  Model condition is the descriptive unit, seeds index stochastic arms, and model families are the resampling clusters for the cross-family quality association \evd{E003}\evd{E015} \\
  \midrule
  Recorded-response intervention &
  The student trains on each 800-sentence outer complement; the untouched 200-sentence block is divided into five inner readout folds &
  Inner-fold transforms, penalty selection, and the fresh response readout &
  Participant-averaged analyses end at five shared outer folds; participant-specific analyses aggregate seed--fold cells within each participant \evd{E004}\evd{E005}\evd{E008}\evd{E011} \\
  \midrule
  Controlled naturalistic voxelwise predictivity assay &
  Five complete-story folds; voxel reliability uses a separate repeated story &
  Training-story \ac{pca}, assembled-design scaling, response centering, one multi-output ridge penalty, and the readout &
  Story-block signs and voxel summaries describe within-participant robustness; they do not create participant inference \evd{E006} \\
  \midrule
  Naturalistic recorded-response intervention &
  Complete stories remain held out from fresh evaluation readouts &
  Training-story transforms, ridge selection, and participant-specific voxel readouts &
  Seeds are technical repetitions within three descriptive participant probes \evd{E013}\evd{E017} \\
  \midrule
  Synthetic-target recovery assay &
  The student trains on 95,999 WikiText sentences and is evaluated on 1,999 fixed held-out sentences &
  The fresh target-readout penalty is selected using training data &
  Equal-coordinate target \(R^2\) is compared within six matched technical seeds \evd{E016} \\
  \midrule
  Saved-student biological-transfer assay &
  Frozen students enter five contiguous Tuckute 800/200 response-readout folds &
  All representation transforms, nuisance transforms, penalties, and fresh response readouts use the 800-sentence complement &
  Six seeds and five folds are aggregated within each of nine participants; participants are the biological units \evd{E025} \\
  \midrule
50-component linear target-projection assay &
  Five contiguous Tuckute folds; aligned targets and the frozen row twin share the train-fold target-\ac{pca} basis &
  Each aligned or twin multi-output readout selects one scalar penalty on its own training complement &
  Folds are aggregated within each of nine participants; there is no student-seed axis \evd{E030} \\
  \midrule
  Saved-student target-retention analysis &
  Frozen students and targets use the same five outer folds &
  Representation and target transforms and the multi-output target readout remain train-fold-local &
  Sums of squares are pooled across target coordinates; paired contrasts end at six technical seeds \evd{E030} \\
  \midrule
  Composed-path diagnostic &
  Both linear paths are fitted on the same outer-training rows and scored once on the untouched test block &
  Target-to-response and student-to-target coefficient blocks are fitted separately inside each fold &
  Folds and six seeds are aggregated within each of nine participants \evd{E030} \\
  \midrule
Practical-utility evaluation &
  Fixed in-domain and out-of-domain corpora are scored without fitting a biological readout &
  No evaluation-time model fitting; corpus-level losses are computed under the fixed scoring protocol &
  Paired arm contrasts end at eight technical seeds and remain conditional on the fixed models and corpora \evd{E009} \\
\end{longtable}
\endgroup

The aggregation order can be written without treating technical repetitions as new scientific units.
Let \(d_{u,r}\) be a matched contrast for final unit \(u\) and technical cell \(r\), and let \(\mathcal A_u\) denote the declared within-unit aggregation:
\begin{equation}
  \widehat\theta_u
  =
  \mathcal A_u\!\left(\{d_{u,r}:r\in\mathcal R_u\}\right),
  \qquad
  \widehat\Theta
  =
  \mathcal I\!\left(\widehat\theta_1,\ldots,\widehat\theta_U\right).
  \label{eq:aggregation-order}
\end{equation}
Here \(\mathcal A_u\) may be a mean, a median, or pooled sums of squares, while \(\mathcal I\) is the descriptive summary or inferential procedure applied to the \(U\) final units.
Equation~\ref{eq:aggregation-order} says that seeds and folds are first reduced within a participant when the participant is the scientific unit.
Analyzing every seed--fold cell as if it were an independent participant would change both the estimator and its uncertainty.

The final unit determines what can generalize.
Seeds vary optimizer randomness under one fixed procedure.
Stimulus folds share observations and have overlapping training complements.
Voxels and \acp{roi} share a participant, preprocessing, and stimuli.
Target coordinates share one generated target and student.
None can replace a participant when the claim concerns people \autocite{lazic2010}.
Conversely, several participants evaluated on one fixed sentence set do not by themselves establish generalization to a population of stimuli.
Shared-fold and story-block summaries are therefore robustness checks, not substitutes for participant or stimulus-population inference. \evd{E006}\evd{E008}\evd{E025}\evd{E030}

\subsection{Uncertainty, multiplicity, and decision rules}

Four choices must remain separate.
The inference unit identifies the observations that enter uncertainty.
The interval or test describes variation across those units.
The multiplicity family identifies which simultaneous claims require adjustment.
The decision rule states what evidence was required to continue to a later stage.
Unless stated otherwise, the reported intervals are two-sided 95\% intervals.

For \(U\) final-unit estimates \(\widehat\theta_1,\ldots,\widehat\theta_U\), the ordinary \(t\) interval is
\begin{equation}
  \overline{\theta}
  \ \pm\
  t_{1-\alpha/2,U-1}
  \frac{s_\theta}{\sqrt U},
  \qquad
  \overline{\theta}=\frac{1}{U}\sum_{u=1}^{U}\widehat\theta_u.
  \label{eq:t-interval}
\end{equation}
Here \(\alpha\) is the two-sided error rate.
Equation~\ref{eq:t-interval} uses the number of final units \(U\), not the number of coordinates, folds, or technical runs.
Its sampling interpretation requires independent final-unit estimates and approximate normality when \(U\) is small.
Participant intervals are conditional on the fixed stimuli, trained students, and analysis protocol used to construct each participant estimate.

For \(K\) matched technical contrasts \(d_1,\ldots,d_K\), the exact two-sided magnitude-preserving sign-flip test used in the seed-level analyses is
\begin{equation}
  p_{\mathrm{flip}}
  =
  \frac{1}{2^K}
  \sum_{\epsilon\in\{-1,+1\}^K}
  \mathbf 1\!\left[
    \left|\frac{1}{K}\sum_{k=1}^{K}\epsilon_kd_k\right|
    \ge
    \left|\frac{1}{K}\sum_{k=1}^{K}d_k\right|
  \right].
  \label{eq:sign-flip}
\end{equation}
This test requires the contrast signs to be exchangeable under the null conditional on their magnitudes.
It differs from the participant sign test, which uses only the number of positive and negative participant effects and requires independent participant signs.
The Wilcoxon signed-rank sensitivity additionally assumes independent and approximately symmetric paired effects.
The shared-fold cluster bootstrap resamples the five outer-fold summaries as clusters.
Because the folds reuse observations through overlapping training complements, both that bootstrap and the fold \(t\) interval are fixed-block stability summaries rather than stimulus-population intervals.
The cross-family quality bootstrap instead resamples model families rather than individual models.

Table~\ref{tab:inference-map} maps each scientific role to its uncertainty procedure and stage-specific decision scope.

\begingroup
\footnotesize
\setlength{\tabcolsep}{3pt}
\renewcommand{\arraystretch}{1}
\begin{longtable}{@{}>{\raggedright\arraybackslash}p{27mm}>{\raggedright\arraybackslash}p{48mm}>{\raggedright\arraybackslash}p{59mm}@{}}
  \caption{Uncertainty procedures and decision rules by scientific role.}
  \label{tab:inference-map}\\
  \toprule
  Scientific role & Interval or test and its unit & Multiplicity and decision scope \\
  \midrule
  \endfirsthead
  \multicolumn{3}{c}{\tablename~\thetable{} (continued)} \\
  \toprule
  Scientific role & Interval or test and its unit & Multiplicity and decision scope \\
  \midrule
  \endhead
  \midrule
  \multicolumn{3}{r}{Continued on next page} \\
  \endfoot
  \bottomrule
  \endlastfoot
  Controlled regional predictivity assay &
  Held-out-fold dispersion and trained--untrained contrasts on the fixed sentence set &
  Prespecified layer and nuisance grids bound the assay; folds do not support participant or stimulus-population inference \evd{E002} \\
  \midrule
  Controlled naturalistic voxelwise predictivity assay &
  Within-participant trained--untrained point contrasts and five story-block signs; no voxel-resampling interval &
  Spatial voxels are not independent units, and overlapping story folds remain robustness checks \evd{E006} \\
  \midrule
Distillation-headroom analysis &
  Fold--seed bootstrap conditional on convergence; cross-family quality association uses a family-cluster bootstrap, with pooled Fisher intervals only as a sensitivity &
  The declared retention bands classify opportunity within this assay and are not universal alignment thresholds \evd{E003}\evd{E015} \\
  \midrule
  Participant-averaged recorded-response intervention &
  Seeds are averaged within each outer fold before a five-fold \(t\) interval &
  The interval describes stability across the five shared blocks of the fixed sentence set \evd{E004}\evd{E005} \\
  \midrule
  Participant-specific recorded-response intervention &
  Participant medians over seed--fold cells receive a participant \(t\) interval and exact one-sided sign test; a shared-fold interval and cluster bootstrap assess block sensitivity &
  Requiring agreement between participant and shared-fold readings is a conservative conjunctive rule, not a joint crossed-population confidence interval \evd{E008}\evd{E011} \\
  \midrule
  Naturalistic recorded-response intervention &
  Seed summaries are technical repetitions within three participant-specific probes &
  The small participant set supports descriptive mechanism probing, not population inference \evd{E013}\evd{E017} \\
  \midrule
  Synthetic-target recovery assay and retained-student movement audit &
  Six matched-seed contrasts receive descriptive intervals and declared exact sign-flip tests &
  Manipulation and language-quality criteria qualify later interpretation but provide no biological replication \evd{E016}\evd{E025} \\
  \midrule
  Comparator-adequacy audit and attribution assessment &
  Fixed-sample margins and paired arm contrasts are reported without coordinate-level \(p\)-values &
  Comparator adequacy is a pre-outcome or baseline design condition for attribution; realized movement is diagnostic, and statistical precision cannot repair a non-identifying comparator \evd{E026}\evd{E030} \\
  \midrule
  Saved-student biological-transfer assay &
  Nine participant effects receive a \(t\) interval, a one-sided upper bound, an exact sign test, and a Wilcoxon sensitivity &
  Seeds and blocks remain descriptive; the upper bound is compared only with the prospectively inherited continuation rule \evd{E025} \\
  \midrule
  50-component linear target-projection assay &
  Participant effects receive \(t\) intervals and exact sign tests &
  Its participant quantities belong to a separate Bonferroni family with a simultaneous one-sided bound \evd{E030} \\
  \midrule
  Saved-student target-retention analysis &
  Six paired technical-seed contrasts receive a descriptive interval and an exact \(2^6\) sign-flip test &
  Target retention has its own adjusted interval family and does not create biological evidence \evd{E030} \\
  \midrule
  Composed-path diagnostic &
  Participant path-product quantities receive \(t\) intervals and exact sign tests &
  The composed path has its own adjusted interval family; it is predictive and is not a causal mediation effect \evd{E030} \\
  \midrule
  Practical-utility evaluation &
  Eight paired technical-seed endpoint observations were recorded, but the completed analysis reports only mean contrasts and an observed-variance sensitivity calculation &
  No paired endpoint interval or equivalence test was retained; the \ac{mde} supplies no hypothesis or equivalence decision \evd{E009} \\
\end{longtable}
\endgroup

For the primary retained-student movement audit, let \(H_s\) and \(H_s^{\mathrm{KD}}\) be layer-6 representation matrices for the same 1,999 untouched WikiText sentences in technical seed \(s\), and let \(C=I-\mathbf{1}\mathbf{1}^{\top}/n\) center examples.
The centered relative-Frobenius distance and linear-\ac{cka} distance are
\begin{equation}
D_{F,s}
=
\frac{\lVert C(H_s-H_s^{\mathrm{KD}})\rVert_F}
{\lVert C H_s^{\mathrm{KD}}\rVert_F},
\qquad
D_{\mathrm{CKA},s}
=
1-
\frac{\lVert (CH_s)^{\top}(CH_s^{\mathrm{KD}})\rVert_F^2}
{\lVert (CH_s)^{\top}(CH_s)\rVert_F\,
 \lVert (CH_s^{\mathrm{KD}})^{\top}(CH_s^{\mathrm{KD}})\rVert_F}.
\label{eq:movement-distances}
\end{equation}
The first quantity measures centered activation change relative to the ordinary-\ac{kd} activation norm.
The second measures the change in centered linear representation geometry, with zero denoting identical linear-\ac{cka} geometry.
Global parameter displacement is reported separately as
\begin{equation}
D_{\theta,s}
=
\frac{\lVert\theta_s-\theta_s^{\mathrm{KD}}\rVert_2}
{\lVert\theta_s^{\mathrm{KD}}\rVert_2}.
\label{eq:parameter-displacement}
\end{equation}
These post-training quantities diagnose the realized intervention; they are not comparator-identification conditions.\evd{E025}\evd{E026}

An \ac{mde} describes the resolution of a planned contrast under its declared variance model, sample size, error rate, and power.
It does not show that a smaller effect is zero or that a larger effect matters biologically.
Likewise, language-quality criteria decide whether two arms are sufficiently comparable for a specified analysis; they do not establish a biological endpoint.

The \result{e030_practical_threshold} unique-\(R^2\) rule was inherited prospectively from the earlier biological-transfer design and then frozen as the continuation reference for the 50-component linear target-projection assay. \evd{E025}\evd{E030}
It can support a statement about whether an assay-specific simultaneous bound reaches that continuation rule.
It cannot define biological importance across another response scale, nuisance model, layer, cohort, or target.
No retrospective choice of a more favorable analysis can change that scope.

\subsection{Information-theoretic boundaries}

Population partial \(R^2\) has an information-theoretic identity only under restrictive assumptions.
For scalar \(Y\), jointly Gaussian \((X,Y,Z)\), correctly specified population-linear conditional means, and unregularized population residual variances,
\begin{equation}
  I(X;Y\mid Z)
  =
  -\frac{1}{2}\log\!\left(1-R_{\mathrm{partial}}^2\right),
  \qquad
  \Delta\mathcal R_{\mathrm{semi}}^2
  =
  \left(1-\mathcal R_N^2\right)R_{\mathrm{partial}}^2.
  \label{eq:gaussian-cmi-r2}
\end{equation}
Equation~\ref{eq:gaussian-cmi-r2} converts a population residual-variance ratio into conditional mutual information under the stated Gaussian model.
The multivariate analogue uses a log-determinant ratio.
The thesis instead estimates finite-sample, cross-validated semipartial differences from separately regularized ridge fits.
Those differences may be negative and depend on the nuisance design, rank, penalty, and held-out substrate, so they are neither estimates nor lower bounds of \(I(X;Y\mid Z)\) \autocite{cover2006}.

\begin{samepage}
A data-processing claim requires a separate conditional-independence premise.
Let \(S\) be stimulus text, \(B\) a recorded brain response, \(T=g(S)\) the deterministic synthetic brain-response target, and \(H\) the retained student state.
The chain
\[
  B\longrightarrow T\longrightarrow H
  \qquad\Longleftrightarrow\qquad
  B\perp H\mid T
\]
would require all information from \(B\) that reaches \(H\) to pass through \(T\), which would imply \(I(B;H)\leq I(B;T)\).
The executed design does not establish this chain because \(S\) directly supplies both \(T\) and \(H\), and it also drives \(B\).
Conditioning on \(S\) changes the question: because \(T\) is deterministic given \(S\), \(I(T;B\mid S)=0\).
That identity does not show that the target is useless.
It shows that conditioning on the exact stimulus cannot establish additional target information beyond that stimulus or recover the desired unconditioned mediation claim.
Data processing therefore supplies a bound only after the relevant conditional independence is defended; it does not predict the sign of a held-out encoding contrast \autocite{cover2006}. \evd{E016}\evd{E030}
\end{samepage}

Output matching and auxiliary-target fitting impose still weaker constraints on retained representations.
Let \(d_H\) be any declared discrepancy between teacher and student internal states.
Then, without additional identifying assumptions,
\[
  D_{\mathrm{KL}}\!\left(p_T^\tau\middle\|q_S^\tau\right)=0
  \ \not\Rightarrow\
  d_H(H_S,H_T)=0,
  \qquad
  \mathcal L_{\mathrm{target}}\downarrow
  \ \not\Rightarrow\
  \Delta R_{\mathrm{transfer}}^2>0.
\]
The first non-implication holds because different internal parameterizations can realize the same output distribution.
The second holds because a temporary target map can absorb part of the auxiliary objective, and a predictable target can be learned without producing participant-general improvement on recorded responses.
Low output or target loss can establish optimization success under its own objective.
It cannot replace fresh tests of retained-student movement, comparator adequacy, brain-response attribution, biological transfer, or practical utility.
Section~\ref{sec:results} reports those empirical stages without treating any formal analogy as their outcome. \evd{E003}\evd{E016}\evd{E025}\evd{E030}
